%% kapa: exam=1 round=22 year=2011 period=2 form=A subject=appraisal_law subject_ko=부동산관계법규 seq=17 number=57
%% source: Q-Net 감정평가사 기출문제 (cst00309) artlSeq=5036790 file=extracted/1008/5036790/2011년 제22회 감정평가사 1차시험 기출문제/2교시.hwp
%% conversion: kordoc 4.24.51 -> tools/gen1.py (auto); review=proofread-gpt-astra; flags=none
\begin{question}{57}
부동산 가격공시 및 감정평가에 관한 법령상 개별공시지가에 관한 설명으로 옳은 것은?
\begin{choices}
  \item 개별공시지가에 표준지 선정의 착오가 있음을 발견한 때에는 지체 없이 이를 정정하여야 한다.
  \item 표준지로 선정된 토지에 대해서도 별도로 개별공시지가를 결정ㆍ공시하여야 한다.
  \item 개별공시지가를 결정ㆍ공시하기 위하여 개별토지의 가격을 산정하고자 하는 경우에는 둘 이상의 감정평가업자에게 이를 의뢰하여야 한다.
  \item 개별공시지가를 결정ㆍ공시함에 있어 필요하다고 인정되는 때에는 토지소유자 그 밖의 이해관계인의 의견청취를 생략할 수 있다.
  \item 개별공시지가에 대한 이의신청은 그 결정ㆍ공시일부터 30일 이내에 국토해양부장관에게 하여야 한다.
\end{choices}
\end{question}
\answer{1}
